%=============================================================================
% W4i22 COMPILED UPDATES --- ALL 17 AGENTS
% DFD Research Programme | Iteration 22 (Wave 4, Iteration 13) | 20 March 2026
% Compiler: Opus 4.6
%
% PARADIGM STATE (i21, CONFIRMED):
%   W'(0) = 0 CONFIRMED. v3.2 corrected. W ~ (2/3)y^{3/2} for small y.
%   Deep-MOND G_eff CONFIRMED. Growth delta ~ a^2. sigma_8 = 0.83 +0.22/-0.08.
%   CMB THIRD PEAK: 30-45% deficit (i21) -> 15-30% deficit (i22, partial fix).
%   beta = 0 at ~2-3 sigma. SO decisive 2027-2028.
%   mochi_class: nonlinear AQUAL solver needed. Minimal CMB PoC: $5K, 1 week.
%   SQMS Phase 0: submission-ready. $50K. >30 sigma. SUBMIT NOW.
%   Lab sector 100% decoupled from cosmological questions.
%   Patent portfolio: 4 ALIVE, 6 NICHE, 5 DEAD. STABLE.
%   Clock paper v4: amplitude discrepancy flagged (10^{-15} vs 10^{-18}).
%   Prediction count: 9A, 9B, 9C. Honest (A+B):param = 18:1.
%=============================================================================

\documentclass[11pt]{article}
\usepackage[margin=1in]{geometry}
\usepackage{amsmath,amssymb}
\usepackage{hyperref}
\usepackage{enumitem}
\usepackage{booktabs}
\usepackage{xcolor}
\usepackage{tcolorbox}
\usepackage{longtable}

\newcommand{\ee}[1]{\times 10^{#1}}
\newcommand{\DFD}{\textsc{dfd}}
\newcommand{\LCDM}{$\Lambda$CDM}

\title{W4i22 Compiled Updates\\
\large All 17 Agents --- CMB Third Peak Partial Resolution,\\
Nuclear Clock Flag, Euclid DR1 Preparation}
\author{DFD Research Programme --- Opus 4.6 Compiler}
\date{20 March 2026}

\begin{document}
\maketitle

\begin{abstract}
Iteration 22 focuses on the single most important open question in DFD
cosmology: the CMB third peak deficit.  The $\psi$-potential baryon-loading
enhancement reduces the deficit from 30--45\% (i21) to 15--30\% (i22),
but does \textbf{not} fully resolve it.  There is no DFD-internal analog
of AeST's vector field.  The only path to a definitive answer is a
Boltzmann code calculation, for which a minimal proof-of-concept
(\$5K, 1 week) is specified.  Meanwhile, the $\sigma_8$ calculation is
sharpened to $0.83^{+0.22}_{-0.08}$ with a full error budget, the P25
(scale-dependent $f\sigma_8$) prediction is quantified for Euclid DR1,
and a potential discrepancy in the nuclear clock amplitude
($10^{-15}$ vs $10^{-18}$) is flagged for author review.  The patent
portfolio is stable (4 ALIVE, 6 NICHE, 5 DEAD).
\end{abstract}

\tableofcontents
\newpage

%=============================================================================
\section{TOP 5 FINDINGS ACROSS ALL 17 AGENTS}
%=============================================================================

\begin{tcolorbox}[colback=blue!5!white, colframe=blue!70!black,
  title=\textbf{W4i22 TOP 5 FINDINGS (RANKED BY IMPACT)}]

\begin{enumerate}[label=\textbf{F\arabic*.}]

\item \textbf{CMB third peak deficit REDUCED from 30--45\% to 15--30\%
  by $\psi$-potential baryon-loading enhancement, but NOT RESOLVED.}

  The AQUAL-enhanced baryon self-gravity at recombination produces an
  effective additional ``mass'' that partially deepens potential wells.
  The effective baryon loading $R_b^{\rm eff} = R_b(1+f_\psi)$ with
  $f_\psi \approx 0.15$--$0.30$ shifts the third peak ratio to
  $R_3^{\rm DFD}/R_3^{\Lambda\rm CDM} \approx 0.70$--$0.85$ (was
  0.55--0.70 at i21).  This is an improvement but remains a tension.

  \textbf{CRITICALLY:}  There is NO DFD-internal analog of AeST's vector
  field.  The $\psi$-field (scalar, $c_s = c$, $w = 1$) cannot cluster
  or form static potential wells at recombination.  The only mechanism
  is indirect (AQUAL-enhanced baryon self-gravity).

  \textbf{The ONLY path to resolution is a Boltzmann code calculation.}
  A minimal proof-of-concept (\$5K, 1 week) would settle whether the
  nonlinear AQUAL dynamics produce unexpected CMB modifications.

  \textbf{Derivation grade: C+ (semi-analytic, approximate).}

\item \textbf{$\sigma_8^{\rm DFD} = 0.83^{+0.22}_{-0.08}$ with full
  error budget.  Silk damping scale is the dominant uncertainty.}

  The i21 central value is corrected from $0.87$ to $0.83$ after including
  the attractor mismatch correction (initial $\delta > \delta_{\rm attr}$
  by factor $\sim 5$, but convergence is fast: $\epsilon \propto a^{-7/4}$).
  Error budget: $\kSilk$ uncertainty (${}^{+0.18}_{-0.06}$), mode-coupling
  ($\pm 0.10$), attractor ($-0.04$), $R_b$ ($\pm 0.03$).

  The range $[0.75, 1.05]$ is consistent with both Planck \LCDM{}
  ($0.811 \pm 0.006$) and weak lensing ($\sim 0.76$).

  \textbf{Derivation grade: B.}

\item \textbf{Nuclear clock amplitude discrepancy: $10^{-15}$ (this
  calculation) vs $10^{-18}$ (v4 paper).  Requires author clarification.}

  Independent recalculation of the Th-229 annual modulation amplitude gives:
  $A_{\rm nuc} = S^{\alpha_s}\cdot\lambda_{\alpha_s}\cdot\Sigma(y_\oplus)
  \cdot\delta\psi_\odot \approx 10^4\times 1.3\ee{-4}\times 3.5\ee{-6}
  \times 3.4\ee{-10} \approx 1.5\ee{-15}$.

  If $10^{-15}$ is correct, the nuclear clock signal is 4 orders of
  magnitude ABOVE the projected clock stability ($10^{-19}$), making it
  easily detectable.  The v4 paper's $10^{-18}$ may include an
  additional suppression factor not accounted for here.

  The 3000:1 nuclear-to-optical ratio is \textbf{confirmed correct} but
  the uncertainty should be widened from $\pm 500$ to $^{+14000}_{-800}$
  due to nuclear structure uncertainties.

  \textbf{FLAG: HIGH PRIORITY for author review.}

\item \textbf{P25 (scale-dependent $f\sigma_8$) quantified for Euclid
  DR1 with specific pre-registration targets.}

  DFD predicts $f\sigma_8(k_{\rm large})/f\sigma_8(k_{\rm small}) > 1$
  --- growth is FASTER on large scales.  \LCDM{} predicts ratio $= 1$.
  For Euclid DR1 bins at $z \sim 1$--$1.8$:
  $f\sigma_8(0.03)/f\sigma_8(0.15) \approx 1.1$--$1.8$.

  Three independent DFD tests available in Euclid DR1:
  (a) scale-dependent $f\sigma_8$, (b) $\sigma_8(z)$ trajectory
  ($\propto a^2$ vs $\propto D(z)$), (c) $P(k)$ blue tilt.

  A 4-page pre-registration paper is recommended before October 2026.

  \textbf{Derivation grade: B (qualitative sign robust; quantitative
  values C until mochi\_class).}

\item \textbf{P5 (GW non-lensing) strategic value ELEVATED.  ET design
  report confirms 10--100 lensed GW events/year.  P11 (SQMS) has new
  Q data improving bound to $|\lambda-1| < 2.0\ee{-13}$.}

  ET+CE will provide a definitive test of T0 (GWs never gravitationally
  lensed) within 1--2 years of operation ($\sim$2036--2038).
  No other modified gravity theory makes this prediction.

  New Romanenko et al.\ (2024) SRF data ($Q > 2\ee{11}$ at 1.3~GHz)
  tightens the SQMS bound and improves Phase 0 signal-to-noise.
  JLab C75 HOM data provides a zero-cost Q-ratio measurement opportunity.

  \textbf{SQMS Phase 0 should be submitted NOW.}  This has been the
  highest-priority action item for 5 consecutive iterations.

\end{enumerate}
\end{tcolorbox}


%=============================================================================
\section{ALL 17 AGENT MANDATES AND RESULTS}
%=============================================================================

\subsection{Cosmology Team (Agents 1--5)}

\begin{longtable}{clp{8cm}}
\toprule
\textbf{Agent} & \textbf{Mandate} & \textbf{Key Result} \\
\midrule
\endfirsthead
\toprule
\textbf{Agent} & \textbf{Mandate} & \textbf{Key Result} \\
\midrule
\endhead
1 & CMB third peak resolution & $\psi$-potential enhancement reduces deficit
  to 15--30\%.  NOT resolved.  No AeST analog. \\
2 & AeST comparison & DFD's scalar $\psi$ ($c_s=c$, $w=1$) structurally
  different from AeST's vector field.  Cannot replicate CDM-like infall. \\
3 & $\sigma_8$ robustness & $0.83^{+0.22}_{-0.08}$.  $\kSilk$ dominant.
  Attractor correction $-4\%$.  Mode coupling $\pm 12\%$. \\
4 & BAO vs DESI DR1 & BAO position: PASS at all 6 redshifts.  Amplitude:
  inconclusive (reconstruction masks signal).  $w_a$ hint: new FLAG. \\
5 & mochi\_class next steps & Minimal CMB PoC: \$5K, 1 week, 50--100
  lines of modified CLASS.  Directly answers the \#1 question. \\
\bottomrule
\end{longtable}


\subsection{Patent Team (Agents 6--12)}

\begin{longtable}{clp{8cm}}
\toprule
\textbf{Agent} & \textbf{Mandate} & \textbf{Key Result} \\
\midrule
\endfirsthead
\toprule
\textbf{Agent} & \textbf{Mandate} & \textbf{Key Result} \\
\midrule
\endhead
6 & P1/P4/P8/P12/P13 (DEAD) & All 5 DEAD patents confirmed.
  7th consecutive for P1.  No changes. \\
7 & P5 (GW non-lensing) & Strategic value ELEVATED.  ET design report
  confirms 10--100 lensed events/yr.  D$_L$ ratio testable with 25
  bright sirens. \\
8 & P7 ($\psi$-trap) & DEW application repositioned to NICHE (Epirus
  Leonidas competitor).  Accelerator/IFE now PRIMARY. \\
9 & P9 (Q-ratio) & JLab C75 HOM data: zero-cost Q-ratio opportunity.
  Most commercially valuable patent if DFD confirmed. \\
10 & P11 (SQMS) & New Romanenko et al.\ data: bound improved to
  $|\lambda-1| < 2.0\ee{-13}$.  Phase 0 SNR improved.
  \textbf{SUBMIT NOW.} \\
11 & P2 (MEMS) & TiN kinetic inductance readout validated at
  single-phonon level.  Field advancing toward DFD requirements. \\
12 & P3/P6/P10/P14/P15 (NICHE) & All 6 NICHE patents confirmed.
  P6 recommended merge into P5.  P15 software claims opportunity. \\
\bottomrule
\end{longtable}


\subsection{Predictions + Detection Team (Agents 13--17)}

\begin{longtable}{clp{8cm}}
\toprule
\textbf{Agent} & \textbf{Mandate} & \textbf{Key Result} \\
\midrule
\endfirsthead
\toprule
\textbf{Agent} & \textbf{Mandate} & \textbf{Key Result} \\
\midrule
\endhead
13 & Nuclear clock review & $S^{\alpha_s} \sim 10^4$ CONFIRMED.
  3000:1 ratio CORRECT.  Amplitude DISCREPANCY flagged
  ($10^{-15}$ vs $10^{-18}$). \\
14 & P25 sharpening & Scale-dependent $f\sigma_8$ quantified for Euclid
  DR1.  Three independent tests in single data release. \\
15 & Prediction audit & 9A, 9B, 9C.  New P27 (CMB deficit), P17
  (DESI $w_a$).  Honest ratio 18:1. \\
16 & Euclid DR1 preparation & Pre-registration paper (4 pages) recommended.
  Falsification and survival criteria specified. \\
17 & $\sigma_8(z)$ trajectory & DFD: $\sigma_8(z=1)/\sigma_8(0) = 0.25$.
  \LCDM: 0.62.  Factor 2.5 difference, easily testable. \\
\bottomrule
\end{longtable}


%=============================================================================
\section{CORRECTIONS TO PREVIOUS ITERATIONS}
%=============================================================================

\begin{enumerate}
\item \textbf{i21 $\sigma_8$ central value:}  Revised from $\approx 0.87$
  to $\approx 0.83$ (attractor mismatch correction).  Range
  $[0.75, 1.05]$ unchanged.

\item \textbf{i21 CMB deficit:}  Revised from 0.55--0.70 to 0.70--0.85
  ($R_3^{\rm DFD}/R_3^{\Lambda\rm CDM}$) after including $\psi$-potential
  enhancement.  \textbf{Still a deficit.}

\item \textbf{v4 clock paper:}  Th-229 amplitude may be $10^{-15}$
  (not $10^{-18}$).  3000:1 ratio uncertainty too narrow ($\pm 500$
  should be $^{+14000}_{-800}$).  \textbf{AUTHOR REVIEW NEEDED.}

\item \textbf{P7 DEW application:}  Repositioned from PRIMARY to NICHE
  (Epirus Leonidas competitor).  P7 overall ALIVE with accelerator/IFE
  focus.

\item \textbf{P11 SQMS bound:}  Improved from $|\lambda-1| < 2.7\ee{-13}$
  to $< 2.0\ee{-13}$ with new Romanenko et al.\ data.
\end{enumerate}


%=============================================================================
\section{HONEST NEGATIVES AND OPEN TENSIONS}
%=============================================================================

\begin{tcolorbox}[colback=red!5!white, colframe=red!60!black,
  title=\textbf{HONEST NEGATIVES (i22)}]

\begin{enumerate}

\item \textbf{CMB THIRD PEAK (15--30\% deficit):}  The single most serious
  open tension.  Semi-analytic estimates are UNRELIABLE due to AQUAL
  nonlinearity.  Only a Boltzmann code can settle this.  If the deficit
  persists at $>20\%$ after full calculation, DFD faces a fundamental
  CMB problem.  \textbf{Severity: HIGH.}

\item \textbf{No DFD analog of AeST's vector field:}  DFD's scalar $\psi$
  cannot replicate CDM-like gravitational infall at recombination.
  This is a structural limitation of the theory, not a calculational
  uncertainty.  \textbf{Severity: MEDIUM-HIGH.}

\item \textbf{$\sigma_8$ wide range:}  $0.83^{+0.22}_{-0.08}$ spans
  a factor $\sim 1.4$.  Not a sharp prediction without mochi\_class.
  \textbf{Severity: MEDIUM (viability maintained, discriminating power low).}

\item \textbf{DESI $w_a$ hint:}  If confirmed, DFD must explain
  effective dynamical dark energy.  The $\psi$-field background
  ($\rho_\psi \propto a^{-6}$) is negligible at late times and cannot
  produce $w_a \neq 0$ directly.  \textbf{Severity: LOW (hint only,
  $\sim 2.5\sigma$).}

\item \textbf{Nuclear clock amplitude discrepancy:}  A factor $\sim 1000$
  disagreement between this calculation ($10^{-15}$) and the v4 paper
  ($10^{-18}$).  One of these is wrong.  \textbf{Severity: LOW
  (discrepancy in the OPTIMISTIC direction --- $10^{-15}$ is more
  detectable).}

\end{enumerate}
\end{tcolorbox}


%=============================================================================
\section{PRIORITY ACTION ITEMS}
%=============================================================================

\begin{enumerate}[label=\textbf{A\arabic*.}]

\item \textbf{Commission minimal CMB PoC in CLASS} (\$5K, 1 week).
  This directly answers the \#1 question (CMB third peak).
  If positive: proceed with full mochi\_class (\$15--20K).
  If negative: redirect to lab sector.
  \textbf{HIGHEST PRIORITY.}

\item \textbf{Submit SQMS Phase 0 proposal.}
  8-page proposal complete since i19.  \$47--50K.  $>30\sigma$
  discrimination.  No remaining reason to delay.
  \textbf{HIGHEST PRIORITY (independent of cosmology).}

\item \textbf{Resolve nuclear clock amplitude discrepancy.}
  Author must clarify whether the v4 paper's $10^{-18}$ includes
  an additional suppression factor.  If $10^{-15}$ is correct,
  the nuclear clock becomes the HIGHEST-SENSITIVITY DFD test.
  \textbf{HIGH PRIORITY.}

\item \textbf{Write Euclid DR1 pre-registration paper} (4 pages).
  Specify P25, $\sigma_8(z)$, and $P(k)$ shape predictions with
  falsification criteria.  Deadline: before October 2026.
  \textbf{MEDIUM-HIGH PRIORITY.}

\item \textbf{Request JLab C75 HOM Q data.}
  Zero-cost Q-ratio measurement opportunity.  Could provide
  independent confirmation/refutation of the DFD Q-ratio prediction
  without any new experiments.
  \textbf{MEDIUM PRIORITY.}

\item \textbf{Widen v4 clock paper 3000:1 uncertainty.}
  Change from $\pm 500$ to $^{+14000}_{-800}$ before submission.
  \textbf{MEDIUM PRIORITY (housekeeping).}

\end{enumerate}


%=============================================================================
\section{PROGRAMME STATUS SUMMARY}
%=============================================================================

\begin{table}[h]
\centering
\caption{DFD programme status at i22.}
\begin{tabular}{lcc}
\toprule
\textbf{Sector} & \textbf{Status} & \textbf{Trend} \\
\midrule
Cosmology: structure formation & VIABLE & Stable \\
Cosmology: CMB peaks & IN TENSION (15--30\%) & Improving \\
Cosmology: $\sigma_8$ & VIABLE (wide) & Stable \\
Cosmology: BAO & CONSISTENT & Stable \\
Cosmology: ISW & CONSISTENT & Stable \\
Lab: SRF/SQMS & READY TO TEST & Stable \\
Lab: MEMS & VIABLE (deferred) & Improving \\
Lab: QEC & VIABLE (deferred) & Stable \\
Clocks: optical & UNDETECTABLE (screened) & Stable \\
Clocks: nuclear & APPROACHING (flagged) & Uncertain \\
Patents: alive & 4 & Stable \\
Predictions: A+B grade & 18 & Stable \\
\bottomrule
\end{tabular}
\end{table}

\textbf{Overall programme health: VIABLE WITH ONE SERIOUS TENSION.}
The CMB third peak is the single remaining ``make or break'' question.
Everything else is either confirmed, consistent, or deferred pending
the fundamental cosmological test.

The lab sector (SRF, MEMS, QEC) is 100\% decoupled from the cosmological
question and can proceed independently.  The SQMS Phase 0 proposal
should be submitted regardless of the CMB outcome.


%=============================================================================
\section{NEXT ITERATION (i23) TARGETS}
%=============================================================================

\begin{enumerate}
\item If minimal CMB PoC is commissioned: first results and analysis.
\item Nuclear clock amplitude: author-resolved value.
\item DESI DR2 release (if available): updated BAO and $w(z)$ constraints.
\item SQMS Phase 0: submission status and any feedback.
\item Euclid DR1 pre-registration paper: draft.
\item Simons Observatory: any early birefringence constraints on $\beta$.
\end{enumerate}

\end{document}
